\documentclass[10pt]{article}

\usepackage[margin=1in]{geometry}
\usepackage{amsmath,amssymb}
\usepackage{booktabs}
\usepackage{enumitem}
\usepackage{float}
\usepackage{microtype}
\usepackage[hidelinks]{hyperref}
\usepackage{xcolor}
\usepackage{xspace}

\newcommand{\cpb}{\textsc{ClaimPromoteBench}\xspace}
\newcommand{\allow}{\texttt{ALLOW}}
\newcommand{\hold}{\texttt{HOLD}}
\newcommand{\demote}{\texttt{DEMOTE}}
\newcommand{\escalate}{\texttt{ESCALATE}}

\title{\cpb: Synthetic, Policy-Constrained Evaluation of Evidence-State Transitions}
\author{Anonymous Author(s)}
\date{}

\begin{document}
\maketitle

\begin{abstract}
Producing evidence does not determine which claim-state transition it licenses. We study this intervening decision as synthetic policy conformance: given a claim, typed evidence, current state, target, and fixed author-defined policy, return \allow, \hold, \demote, or \escalate. \cpb uses controlled pairs and counts every service or output failure as incorrect. In a preregistered hidden study of 300 items and 150 pairs, a Qwen configuration returned 272/300 valid decisions; end-to-end item/pair accuracy was .8600/.7533 versus .9485/.9187 conditional accuracy at .9067/.8200 decision/pair coverage. A public, preregistered robustness challenge of 180 items and 90 pairs gave Qwen/Flash end-to-end accuracies of .7333/.5333 and .5000/.1444; conditional values were .8919/.7500 and .6164/.2131. Complete confusion matrices show \demote{} recall of 4/30 and 0/30. A prospectively frozen model-family breadth supplement ran seven service configurations on the same known public challenge: 1,260 terminal cells, zero retries, and 78 valid responses. Route-level end-to-end item/pair accuracy ranged from .0000--.2167/.0000--.0444; conditional metrics were undefined for four routes and used at most 65 valid outputs elsewhere. The old r6 Flash lane records 300 HTTP 401 failures on its old frozen service interface. Because model, serving stack, and adapter differ, these outcomes are descriptive rather than causal or model rankings. The deterministic executor is a definitional policy ceiling; Naive Bayes remains a weak baseline.
\end{abstract}

\section{Introduction}

A valid artifact licenses only the scope it checks. A finite computation supports its registered finite domain, while a proof certificate supports its exact encoded statement in a pinned environment. Independent reproduction changes provenance rather than scope. A system that collapses these distinctions can assign a claim status stronger than the evidence warrants.

We isolate the decision between evidence production and claim-state assignment:
\begin{quote}
\emph{Given a claim, typed evidence, a current state, and a requested transition, which action conforms to a fixed, machine-checkable policy?}
\end{quote}
The benchmark treats this question as synthetic policy conformance under an authored operational definition. Oracle agreement therefore measures behavior relative to that definition, not scientific truth or community consensus. The task excludes open-ended discovery, experimental design, literature synthesis, and general research-agent competence.

\cpb records research status, computational support, formal certification, and independent review on separate state axes. A controlled pair fixes the claim and requested transition while changing one registered prerequisite. We score \allow, \hold, \demote, and \escalate, retaining unusable responses in the end-to-end denominator. Coverage measures whether a service returns a valid decision; conditional accuracy scores valid responses; end-to-end item and pair accuracy score every scheduled case.

\begin{samepage}
Our contributions are:
\begin{enumerate}[leftmargin=*]
  \item an executable controlled-pair protocol that maps typed evidence and requested state transitions to four actions under a fixed authored policy;
  \item failure-aware item and pair estimands, applied in a preregistered hidden diagnostic, a precommitted action-balanced public challenge, and a prospectively frozen model-family breadth supplement; and
  \item a replayable evidence artifact that separates service output from a definitional policy ceiling and produces privacy-filtered reviewer derivatives.
\end{enumerate}
\end{samepage}

\section{Related Work}

\paragraph{Behavioral and contrastive evaluation.}
CheckList organizes behavioral tests around capabilities and expected invariances \cite{ribeiro2020checklist}, while contrast sets probe local decision boundaries through small, meaningful input changes \cite{gardner2020contrastsets}. Metamorphic testing similarly checks relations among transformed inputs when a conventional test oracle is difficult to apply \cite{chan1998metamorphic}. \cpb uses authored policy labels and controlled pairs: most pairs intentionally change the action, while separate metamorphic checks preserve the action under identifier renaming, evidence reordering, and irrelevant-evidence insertion. These tests characterize selected local behaviors; they are not an estimate of natural-distribution performance.

\paragraph{Structured output and selective prediction.}
Incremental parsing and grammar-constrained decoding can prevent structurally invalid generations \cite{scholak2021picard,geng2023grammar}. They cannot guarantee that a structurally valid action satisfies cross-field semantics or the evidence policy. Selective classification studies risk--coverage trade-offs and reject options \cite{geifman2017selective,geifman2019selectivenet}; recent abstention benchmarks expose continuing failures on unanswerable tasks \cite{kirichenko2025abstentionbench}. We therefore report coverage and conditional correctness together, while end-to-end metrics count every absent or invalid decision as incorrect. \hold{} is a policy action about current evidence, whereas service failure is recorded as no valid decision.

\paragraph{Explicit rule following and formal specifications.}
RuleArena evaluates long-context natural-language rule application for airline baggage fees, NBA transactions, and tax regulations, including logical and numerical reasoning \cite{zhou2025rulearena}. OSVBench instead asks models to generate formal specifications for operating-system verification \cite{li2026osvbench}. \cpb neither answers domain questions nor generates specifications: it applies a complete machine-readable evidence policy to a typed ledger and scores a four-way state-transition action. Its oracle is an authored synthetic definition, not an external judgment of scientific truth.

\paragraph{Scientific evidence and research-agent evaluation.}
CLAIM-BENCH, DeepFact, and ClaimDB evaluate scientific claim extraction or factual verification \cite{javaji2025claimbench,huang2026deepfact,theologitis2026claimdb}. ScienceAgentBench, CORE-Bench, and POPPER evaluate executable research tasks, computational reproduction, or hypothesis falsification \cite{chen2025scienceagentbench,siegel2025corebench,huang2025popper}. SeekBench, Omni-Decision, EviGraph, and ClaimLedger emphasize grounding or explicit evidence maintenance \cite{shao2025seekbench,ma2026omnidecision,ren2026evigraph,elliman2026claimledger}. Unlike these broader tasks, \cpb fixes an authored policy and scores only the status-transition decision after evidence registration.

\paragraph{Evidence ledgers and traceability.}
Evidence-Ledger Adjudication assigns externally labeled support relations and routes claims on a 2,335-row benchmark \cite{chen2026evidenceledger}. LEDGER reconstructs claim-support paths and artifact lineage from observed agent sessions \cite{kim2026ledger}. \cpb instead tests executable conformance of state transitions to an authored synthetic policy through controlled pairs; it neither predicts external support labels nor reconstructs session traces.

\section{Task and Evidence Flow}

Each item is \(x=(c,E,s,\tau,t)\): a claim \(c\), typed evidence \(E\), current state \(s\), requested target \(\tau\), and frozen evaluation time \(t\). The state has four axes:
\begin{align*}
\texttt{research\_status} &\in \{\texttt{PROPOSED},\texttt{OPEN},\texttt{RESOLVED}\},\\
\texttt{computation\_status} &\in \{\texttt{NONE},\texttt{COMPUTATIONAL}\},\\
\texttt{formal\_status} &\in \{\texttt{NONE},\texttt{FORMALLY\_CERTIFIED}\},\\
\texttt{review\_status} &\in \{\texttt{UNREVIEWED},\texttt{INDEPENDENTLY\_VERIFIED}\}.
\end{align*}
The public policy checks dependencies, scope, freshness, provenance, runtime completion, contradictions, statement identity, and artifact integrity. Its deterministic oracle \(\pi_P(x)\) evaluates the following rows from top to bottom; the first matching row fixes both the action and proposed state.

\begin{table}[H]
\centering
\footnotesize
\begin{tabular}{@{}p{0.12\linewidth}p{0.62\linewidth}p{0.18\linewidth}@{}}
\toprule
Action & First-match condition & Proposed state \\
\midrule
\demote{} & A currently elevated axis has an absent, broken, stale, mismatched, or policy-nonconforming recorded basis. & Lower each affected axis. \\
\escalate{} & A fresh, identity-matched, independently produced contradiction is open and passes its runtime and verifier checks. & Keep current state. \\
\hold{} & The requested promotion lacks a complete valid dependency closure or any required scope, freshness, runtime, verifier, provenance, or axis-specific check. & Keep current state. \\
\allow{} & None of the preceding conditions holds. & Accept target state. \\
\bottomrule
\end{tabular}
\end{table}

The model receives the complete versioned policy contract, a canonical JSON evidence ledger, and passive verifier checks that report raw facts but no transition decision. Action precedence is \demote{} \(>\) \escalate{} \(>\) \hold{} \(>\) \allow; reason codes follow the policy's published total order. The response must be one JSON object with exactly six fields: \texttt{action}, \texttt{proposed\_state}, \texttt{p\_admissible}, \texttt{used\_evidence\_ids}, \texttt{reason\_codes}, and \texttt{missing\_evidence\_ids}. No rationale text is requested.

Let \(v(x)\) indicate a syntactically, schema-, and semantically valid response and let \(\hat y(x)\) be its action. End-to-end item correctness is
\[
C_{\mathrm{item}}(x)=\mathbb{1}[v(x)=1]\mathbb{1}[\hat y(x)=\pi_P(x)].
\]
For controlled pair \(i=(x_i^a,x_i^b)\), \(C_{\mathrm{pair}}(i)=C_{\mathrm{item}}(x_i^a)C_{\mathrm{item}}(x_i^b)\). Decision coverage is the fraction with \(v=1\); pair coverage is the fraction with both members valid. Conditional accuracies divide correct items or pairs by those coverage denominators. End-to-end accuracies divide by all scheduled items or pairs.

\begin{figure}[H]
\centering
\footnotesize
\setlength{\tabcolsep}{3pt}
\begin{tabular}{@{}cc@{}}
\multicolumn{2}{p{0.92\textwidth}}{\centering\textbf{Fixed context:} arithmetic claim; \texttt{formal\_status=NONE}; target \texttt{FORMALLY\_CERTIFIED}; a passing certificate depends on one registered dependency.}\\[3pt]
\fbox{\parbox{0.41\textwidth}{\textbf{Member A}\\Dependency \texttt{verifier\_result=FAIL}\\Gold: \hold{}\\Reason:\\[-1pt]\texttt{DEPENDENCY\_NOT\_VERIFIED}}} &
\fbox{\parbox{0.41\textwidth}{\textbf{Member B}\\Dependency \texttt{verifier\_result=PASS}\\Gold: \allow{}\\Reason:\\[-1pt]\texttt{POLICY\_SATISFIED}}}\\[2pt]
\multicolumn{2}{c}{\textit{Only controlled change:} \texttt{FAIL} \(\rightarrow\) \texttt{PASS}}
\end{tabular}
\vspace{4pt}

\begin{tabular}{@{}c@{\;}c@{\;}c@{\;}c@{\;}c@{\;}c@{\;}c@{\;}c@{\;}c@{}}
\fbox{\parbox[c][0.42in][c]{0.17\linewidth}{\centering Contract + one\\ledger member}}
&$\rightarrow$&
\fbox{\parbox[c][0.42in][c]{0.12\linewidth}{\centering One service\\call}}
&$\rightarrow$&
\fbox{\parbox[c][0.42in][c]{0.13\linewidth}{\centering Exact six-field\\validation}}
&$\rightarrow$&
\fbox{\parbox[c][0.42in][c]{0.14\linewidth}{\centering Seal record,\\then load gold}}
&$\rightarrow$&
\fbox{\parbox[c][0.42in][c]{0.16\linewidth}{\centering E2E analysis +\\private-safe replay}}
\end{tabular}
\caption{A concrete controlled pair and the evidence flow. The public challenge pair changes only \texttt{/evidence/1/verifier\_result}; fixed identifiers and hashes are omitted from the diagram. End-to-end analysis retains every terminal failure.}
\label{fig:evidence-flow}
\end{figure}

Within a controlled pair, only registered fields relevant to one factor change. The illustrated member changes a dependency check from \texttt{FAIL} to \texttt{PASS}; other mappings send a missing dependency to \hold, a current-state binding failure to \demote, and an active contradiction to \escalate. Separate invariance tests rename identifiers, reorder evidence, or add irrelevant records and require both the oracle and rendered prompt semantics to remain unchanged.

\begin{table}[H]
\centering
\scriptsize
\caption{Factor inventory, grouped by policy family. Labels humanize the machine-readable factor identifiers; counts are 20 for r6 and 15 for the challenge.}
\label{tab:factor-inventory}
\begin{tabular}{@{}p{0.20\linewidth}p{0.36\linewidth}p{0.36\linewidth}@{}}
\toprule
Family & r6 factors & Public-challenge factors \\
\midrule
Dependency & dependency statement, dependency verifier, direct dependency, transitive dependency & dependency verifier \\
Formal identity & claim identity, statement identity, trust base & --- \\
Numerical, current basis & domain, precision, preregistration, repetitions & current computation binding, staleness, verifier; current review environment, identity; repetitions \\
Provenance & cold start, review environment, reviewer identity & reviewer identity \\
Runtime & run completion, tool exit, verifier execution & tool exit \\
Freshness, contradiction & age, contradiction, current binding & age; contradiction claim identity, freshness, reviewer identity, status, verifier execution \\
\bottomrule
\end{tabular}
\end{table}

\section{Benchmark and Evaluation Design}

\subsection{Synthetic benchmark}

Version~1 contains 200 pairs (400 items): 50 public-development pairs and 150 hidden pairs over the 20 factors in Table~\ref{tab:factor-inventory}. Numerical cases recompute registered finite-grid values and scope, while formal cases compile generated Lean sources under a pinned toolchain. These checks establish internal policy consistency, not external scientific validity.

The fresh r6 realization contains 150 hidden pairs and 300 unique prompts from the same ontology and mostly the same template family. A preregistered construction commitment and request schedule preceded inference; all 600 terminal cells for two service lanes were sealed before gold loading. Each item received one call with zero client retries. The gold distribution is imbalanced: 150 \allow, 134 \hold, eight \demote, and eight \escalate{} items. We also froze a four-class multinomial Naive Bayes classifier (\(\alpha=1\)) trained on identifier-normalized unigram and bigram prompt counts from the 100 public-development items.

\subsection{Balanced public robustness challenge}

The separately precommitted public challenge tests action asymmetry over the 15 factors in Table~\ref{tab:factor-inventory}, with six pairs per factor. Its 90 pairs contain 90 \allow{} members and 30 members for each of \hold, \demote, and \escalate. We froze the dataset, labels, prompts, adapters, analysis, and execution contract before inference. Each service route made 180 calls with zero client retries. This public, preregistered robustness challenge tests a known policy and public synthetic construction; it does not estimate hidden-test or natural-distribution generalization.

The two routes differ in model and service stack, including output-constraining behavior. We report descriptive service-configuration outcomes because the design cannot identify the effect of model weights, provider, schema mode, or any single component.

\subsection{Prospectively frozen public model-family breadth}

We subsequently froze seven service configurations on the same known public 180-item challenge: new Qwen and DeepSeek configurations and five additional model families. The 1,260 terminal cells use identical prompt bytes, a 1,536-token prompt-only JSON contract, item-major cyclic lane rotation, one application POST per cell, and zero client retries. The schedule, adapters, and analysis were committed before any call, and no output selected a model or item.

This supplement adds model-family and configuration breadth while holding one policy, ontology, and public challenge fixed. Model identity, serving implementation, and API adapter are not causally isolated. The observed configuration outcomes therefore cannot estimate the effect of any one component or hidden generalization.

\section{Results}

\subsection{Coverage, conditional accuracy, and end-to-end accuracy}

The r6 Qwen configuration returned 272/300 valid decisions and 123/150 complete pairs (Table~\ref{tab:core-results}). Conditional item/pair accuracy was .9485/.9187; retaining all scheduled failures reduced end-to-end accuracy to .8600/.7533. Fourteen HTTP 500 responses, one HTTP 503 response, ten timeouts, and three semantic-invalid outputs remain incorrect. The frozen Naive Bayes comparator scored .6133 item and .2267 pair accuracy. It checks a low-cost learned path but is a weak baseline, not a representative competitive comparator.

\begin{table}[H]
\centering
\scriptsize
\caption{Coverage and accuracy. Conditional values use only valid decisions or complete pairs; end-to-end (E2E) values retain all scheduled failures. ``Old'' identifies the frozen r6 interface. N/E denotes undefined conditional correctness and non-evaluable model efficacy because no model output was returned.}
\label{tab:core-results}
\begin{tabular}{@{}lrrrrrr@{}}
\toprule
Configuration & Valid & E2E item & Cond. item & Pair cov. & E2E pair & Cond. pair \\
\midrule
r6 Qwen3.7-Max & 272/300 & .8600 & .9485 & .8200 & .7533 & .9187 \\
r6 public-dev text NB & 300/300 & .6133 & .6133 & 1.0000 & .2267 & .2267 \\
r6 old Flash interface & 0/300 & .0000 & N/E & .0000 & .0000 & N/E \\
\midrule
Public challenge Qwen & 148/180 & .7333 & .8919 & .7111 & .5333 & .7500 \\
Public challenge Flash & 146/180 & .5000 & .6164 & .6778 & .1444 & .2131 \\
\bottomrule
\end{tabular}
\end{table}

The r6 Flash lane returned 300/300 HTTP 401 authorization failures on the old frozen service interface and no model output. Its scheduled end-to-end scores are 0/300 items and 0/150 pairs; conditional correctness, model efficacy, and the registered second confirmation are non-evaluable. The later challenge route returned model outputs under a different service stack, establishing only that its own path executed. It neither repairs r6 nor supports a strict causal comparison.

On the balanced challenge, Qwen returned 148/180 valid decisions and 64/90 complete pairs; Flash returned 146/180 valid decisions and 61/90 complete pairs. Their end-to-end item/pair accuracies were .7333/.5333 and .5000/.1444, while conditional accuracies were .8919/.7500 and .6164/.2131. Qwen recorded 31 transport failures and one semantic failure; Flash recorded 34 semantic failures. The larger pair-level gaps show why conditional correctness cannot replace coverage or the end-to-end estimand.

\subsection{Model-family breadth on the known public challenge}

The seven frozen configurations returned 78/1,260 valid decisions and 52 correct decisions. Only 25/630 pairs had two valid members, of which seven were correct. Table~\ref{tab:breadth-results} reports each route separately: four routes returned no valid decision, Qwen3.8-Max returned one, Kimi-K3 returned 12, and MiniMax-M3 returned 65. Conditional values of 1.0000 for Qwen3.8-Max and Kimi-K3 rest on one and 12 valid items, respectively. MiniMax-M3 reached .2167/.0333 end-to-end item/pair accuracy; its .6000/.1429 conditional values use 65 valid items and 21 covered pairs.

\begin{table}[H]
\centering
\scriptsize
\caption{Prospectively frozen model-family breadth results. Each provider-catalog label denotes one complete service configuration. N/E denotes an empty conditional denominator.}
\label{tab:breadth-results}
\begin{tabular}{@{}lrrrrrr@{}}
\toprule
Configuration & Valid & E2E item & Cond. item & Pair cov. & E2E pair & Cond. pair \\
\midrule
Qwen3.8-Max & 1/180 & .0056 & 1.0000 & .0000 & .0000 & N/E \\
DeepSeek-V4-Pro & 0/180 & .0000 & N/E & .0000 & .0000 & N/E \\
GLM-5.3 & 0/180 & .0000 & N/E & .0000 & .0000 & N/E \\
Kimi-K3 & 12/180 & .0667 & 1.0000 & .0444 & .0444 & 1.0000 \\
MiMo-V2.5-Pro & 0/180 & .0000 & N/E & .0000 & .0000 & N/E \\
MiniMax-M3 & 65/180 & .2167 & .6000 & .2333 & .0333 & .1429 \\
Hunyuan-HY3 & 0/180 & .0000 & N/E & .0000 & .0000 & N/E \\
\bottomrule
\end{tabular}
\end{table}

The 1,182 invalid cells comprise 630 transport, 312 JSON-syntax, 203 schema, and 37 semantic failures. Their timing and status fields do not identify a cause, so we do not attribute them to model weights, serving infrastructure, subscription state, or adapters. All remain in the end-to-end denominator.

No breadth route returned a correct \demote{} decision (Table~\ref{tab:breadth-action-recall}). Kimi-K3 recovered 8/90 \allow{}, 2/30 \hold{}, and 2/30 \escalate{} cases. MiniMax-M3 recovered 32/90, 3/30, and 4/30, respectively. These supports rule out a useful route ranking: most scheduled cells were invalid, and the conditional estimands select very different subsets.

\begin{table}[H]
\centering
\scriptsize
\caption{Breadth end-to-end action recall. No-valid cells are misses; Macro-F1 uses all four actions with zero division set to zero.}
\label{tab:breadth-action-recall}
\begin{tabular}{@{}lrrrrrr@{}}
\toprule
Configuration & \allow{} & \hold{} & \demote{} & \escalate{} & Bal. recall & Macro-F1 \\
\midrule
Qwen3.8-Max & .0111 & .0000 & .0000 & .0000 & .0028 & .0055 \\
DeepSeek-V4-Pro & .0000 & .0000 & .0000 & .0000 & .0000 & .0000 \\
GLM-5.3 & .0000 & .0000 & .0000 & .0000 & .0000 & .0000 \\
Kimi-K3 & .0889 & .0667 & .0000 & .0667 & .0556 & .1033 \\
MiMo-V2.5-Pro & .0000 & .0000 & .0000 & .0000 & .0000 & .0000 \\
MiniMax-M3 & .3556 & .1000 & .0000 & .1333 & .1472 & .2072 \\
Hunyuan-HY3 & .0000 & .0000 & .0000 & .0000 & .0000 & .0000 \\
\bottomrule
\end{tabular}
\end{table}

Table~\ref{tab:breadth-confusion} gives the complete per-configuration confusion matrices in a compact form. No route predicts \demote. MiniMax-M3 instead maps 15/30 gold-\demote{} items to \allow{} and one to \hold; the other 14 have no valid decision. This mirrors the action asymmetry found in r6 and the balanced challenge, but under substantially lower route coverage.

\begin{table}[H]
\centering
\scriptsize
\setlength{\tabcolsep}{3.5pt}
\caption{Complete breadth confusion matrices. Each gold-action cell is predicted \allow/\hold/\demote/\escalate/no-valid; gold supports are 90/30/30/30.}
\label{tab:breadth-confusion}
\begin{tabular}{@{}lcccc@{}}
\toprule
Configuration & Gold \allow{} & Gold \hold{} & Gold \demote{} & Gold \escalate{} \\
\midrule
Qwen3.8-Max & 1/0/0/0/89 & 0/0/0/0/30 & 0/0/0/0/30 & 0/0/0/0/30 \\
DeepSeek-V4-Pro & 0/0/0/0/90 & 0/0/0/0/30 & 0/0/0/0/30 & 0/0/0/0/30 \\
GLM-5.3 & 0/0/0/0/90 & 0/0/0/0/30 & 0/0/0/0/30 & 0/0/0/0/30 \\
Kimi-K3 & 8/0/0/0/82 & 0/2/0/0/28 & 0/0/0/0/30 & 0/0/0/2/28 \\
MiMo-V2.5-Pro & 0/0/0/0/90 & 0/0/0/0/30 & 0/0/0/0/30 & 0/0/0/0/30 \\
MiniMax-M3 & 32/2/0/4/52 & 2/3/0/0/25 & 15/1/0/0/14 & 0/2/0/4/24 \\
Hunyuan-HY3 & 0/0/0/0/90 & 0/0/0/0/30 & 0/0/0/0/30 & 0/0/0/0/30 \\
\bottomrule
\end{tabular}
\end{table}

\subsection{Factor-level pair outcomes}

Factor-level pair results vary sharply (Table~\ref{tab:factor-results}). R6 Qwen has two zero-accuracy factors and seven perfect factors; its weakest cases are dependency statement (0/7 correct, 2/7 complete) and current binding (0/8 correct, 6/8 complete). Challenge Qwen has three zero-accuracy factors and two perfect factors, while challenge Flash has 11 zero-accuracy factors and no perfect factor. Supports of six to eight pairs make these descriptive breakdowns, not separate confirmatory tests. The reviewer artifact retains every factor-level complete/correct count.

\begin{table}[H]
\centering
\scriptsize
\caption{Factor-level pair diagnostics. Coverage and E2E ranges are minima--maxima across factors; zero/perfect counts use E2E pair accuracy. The reviewer artifact contains every factor row.}
\label{tab:factor-results}
\begin{tabular}{@{}lrrrr@{}}
\toprule
Configuration & Factors & Pair coverage range & E2E pair range & Zero / perfect \\
\midrule
r6 Qwen & 20 & .2857--1.0000 & .0000--1.0000 & 2 / 7 \\
Challenge Qwen & 15 & .5000--1.0000 & .0000--1.0000 & 3 / 2 \\
Challenge Flash & 15 & .0000--1.0000 & .0000--.8333 & 11 / 0 \\
\bottomrule
\end{tabular}
\end{table}

\subsection{Action recall and full confusion matrices}

Table~\ref{tab:action-recall} shows end-to-end recall, so no-valid-decision cells count as misses. R6's eight-item \demote{} and \escalate{} supports are too small for a uniform four-action claim. The balanced challenge supplies 30 members per non-\allow{} action and preserves the same qualitative \demote{} weakness: Qwen recovers 4/30 and Flash 0/30.

\begin{table}[H]
\centering
\scriptsize
\caption{End-to-end recall by gold action; no-valid cells are misses. Gold supports are 150/134/8/8 for r6 and 90/30/30/30 for each challenge lane. Macro-F1 averages all four actions with zero division set to zero.}
\label{tab:action-recall}
\begin{tabular}{@{}lrrrrrr@{}}
\toprule
Configuration & \allow{} & \hold{} & \demote{} & \escalate{} & Bal. recall & Macro-F1 \\
\midrule
r6 Qwen3.7-Max & .9533 & .8209 & .0000 & .6250 & .5998 & .6265 \\
Public challenge Qwen & .8778 & .8667 & .1333 & .7667 & .6611 & .7004 \\
Public challenge Flash & .7556 & .3000 & .0000 & .4333 & .3722 & .3928 \\
\bottomrule
\end{tabular}
\end{table}

Table~\ref{tab:confusion} gives every end-to-end confusion cell. The final column distinguishes absence of a valid decision from any policy action. R6 Qwen maps all eight gold \demote{} items to \allow. On the public challenge, Qwen's 20 valid gold-\demote{} responses split across all four actions, while Flash never predicts \demote. Flash also maps 13 gold-\hold{} items to \allow{} and loses 14/30 gold-\escalate{} items to invalid responses.

\begin{table}[H]
\centering
\scriptsize
\caption{Full confusion matrices. Rows are gold actions; action columns are valid predictions; ``No valid'' includes service, syntax, schema, and semantic failures.}
\label{tab:confusion}
\begin{tabular}{@{}llrrrrr@{}}
\toprule
Configuration & Gold & \allow{} & \hold{} & \demote{} & \escalate{} & No valid \\
\midrule
r6 Qwen & \allow{} & 143 & 0 & 0 & 0 & 7 \\
& \hold{} & 4 & 110 & 0 & 2 & 18 \\
& \demote{} & 8 & 0 & 0 & 0 & 0 \\
& \escalate{} & 0 & 0 & 0 & 5 & 3 \\
\midrule
Challenge Qwen & \allow{} & 79 & 0 & 0 & 0 & 11 \\
& \hold{} & 0 & 26 & 0 & 0 & 4 \\
& \demote{} & 8 & 4 & 4 & 4 & 10 \\
& \escalate{} & 0 & 0 & 0 & 23 & 7 \\
\midrule
Challenge Flash & \allow{} & 68 & 4 & 0 & 8 & 10 \\
& \hold{} & 13 & 9 & 0 & 0 & 8 \\
& \demote{} & 25 & 3 & 0 & 0 & 2 \\
& \escalate{} & 0 & 3 & 0 & 13 & 14 \\
\bottomrule
\end{tabular}
\end{table}

\subsection{Definitional policy ceiling}

The deterministic policy executor reads the same model-visible item and public policy that define the labels. It reproduces 300/300 r6 items and 150/150 pairs, all 360 original public-challenge item cells and 180 pair cells, and all 1,260 breadth cells and 630 pairs. This perfect agreement is true by definition of the evaluation target. We call it a \emph{definitional policy ceiling}; it is neither a learned baseline nor independent evidence that the authored policy is scientifically correct. Post-generation gates driven by the same executor are implementation checks and provide no evidence of learned model competence.

\section{Discussion}

\paragraph{What r6 supports.}
The preregistered r6 diagnostic supports one time-stamped Qwen service configuration on a fresh hidden realization of the same synthetic ontology. It handled 258/300 items and 113/150 pairs correctly end to end, with .9067 decision coverage and .8200 pair coverage. Its .9485/.9187 conditional item/pair accuracy answers a narrower question: correctness after a usable response exists. Neither estimand measures broad research ability or naturalistic scientific judgment.

\paragraph{Action balance changes the picture.}
R6 is dominated by \allow{} and \hold, and its 0/8 \demote{} recall precludes a uniform four-action claim. The balanced challenge gives 30 cases to each non-\allow{} action; Qwen recovers 4/30 \demote{} cases and Flash recovers none. Every configuration also has a zero-accuracy factor. The confusion matrices distinguish the remaining errors: Qwen mixes valid demotions across all four actions, while Flash often substitutes \allow{} for \hold{} or \demote. Action recall, factor breakdowns, and pair accuracy expose failures that aggregate item accuracy hides.

\paragraph{Breadth exposes route failure rather than a model ranking.}
The prospectively frozen supplement returns only 78/1,260 valid decisions across seven configurations. Four routes have undefined conditional accuracy; Qwen3.8-Max and Kimi-K3 have conditional item accuracy of 1.0000 on one and 12 valid responses. MiniMax-M3 has the largest valid subset, yet reaches only .2167/.0333 end-to-end item/pair accuracy and .6000/.1429 conditionally. Its complete matrix includes 15 gold-\demote{} items mapped to \allow{} and none mapped correctly. Across all seven routes, \demote{} recall is zero and balanced recall ranges from .0000 to .1472. These values characterize frozen service configurations on one known public challenge; they do not order the underlying model families.

\paragraph{Service reliability is part of end-to-end behavior.}
Every transport, syntax, schema, and semantic failure is a failed scheduled decision. Conditional reporting separates response validity from policy correctness, but it cannot replace the end-to-end estimand. The breadth run records 630 transport failures and 552 output failures; neither status nor timing identifies their cause. The old r6 Flash lane documents 300/300 authorization failures on its old frozen service interface and no model output. The later Flash route shows only that a different path executed. Because the model, service stack, and adapter changed, the Qwen--Flash gap is descriptive and is not a causal ablation.

\paragraph{Comparators and enforcement.}
The public-development Naive Bayes model gives a reproducible low-cost reference, but its .6133 item accuracy, .2267 pair accuracy, .3284 balanced recall, and .3124 macro-F1 make it a weak baseline. Stronger learned comparators remain needed. The deterministic executor reaches perfect agreement because it defines the gold actions. It is a definitional policy ceiling, and treating it as learned model evidence would be circular.

\section{Limitations}

The benchmark is synthetic and policy-constrained, and the experiments demonstrate the protocol on selected service configurations rather than establish broad benchmark validity. R6 reduces post hoc adaptation through a fresh hidden realization, but it retains the same ontology and mostly the same template family. The public, precommitted challenge tests an action-balanced construction rather than hidden-test generalization. Its 90 pairs and 15 factors broaden action coverage, although six pairs per factor cannot support precise factor-specific estimates. The full factor matrix is diagnostic: zero cells identify failures but do not establish stable factor rankings.

Coverage remains a material limitation. R6 loses 28/300 decisions and 27/150 complete pairs; the public challenge loses 32/180 Qwen decisions and 34/180 Flash decisions. Conditional accuracies of .9485/.9187, .8919/.7500, and .6164/.2131 must therefore be read beside their end-to-end counterparts. Thirty challenge cases per non-\allow{} action still support only coarse recall estimates, including the observed 4/30 and 0/30 \demote{} results.

The breadth supplement is more severely coverage-limited: 1,182/1,260 cells are invalid, four routes have no valid output, and only 25/630 pairs are covered. Its route-level end-to-end item/pair ranges of .0000--.2167/.0000--.0444 are well-defined, but most conditional cells are undefined or use at most 65 items and 21 pairs. The complete per-configuration confusion matrices preserve all failures and show zero \demote{} recall throughout. They do not support precise action comparisons, stable family rankings, or an inference that a zero-valid route lacks the underlying policy skill.

The evidence covers a few time-stamped service configurations and one call per item. Zero retries define the estimand but do not characterize retry policies, sampling variance, or fixed model weights. The breadth supplement increases the configuration count but keeps one public challenge, policy, ontology, prompt contract, and subscription endpoint. Its adapters and serving implementations can differ, so a route that returns no valid decision under the frozen contract does not estimate model-family competence. The old and new Flash routes likewise cannot identify why the former returned 401 or what its r6 model accuracy would have been. No strict causal provider, model, or structured-output comparison follows.

The policy and labels are authored operational definitions. There is no human annotation study, external policy consensus, agreement measure, or validation against downstream scientific outcomes. Before any \cpb call, we froze and reviewed 12 pointer-only CORE-Bench candidates; none contained the time-stamped artifacts, dependency chain, runtime status, and provenance needed for a transition. We therefore created zero naturalistic benchmark cases and made zero model calls. No human annotators or agreement statistic are claimed. Naturalistic transfer remains open, as do stronger learned baselines, independent policy assessment, and downstream utility.

\section{Reproducibility and Artifact Policy}

The repository binds policies, generators, datasets, prompts, terminal records, gold-loading order, analyses, and receipts. R6 freezes its construction and 600-request schedule before inference and seals every terminal record before gold loading. The public challenge freezes its dataset, action balance, prompt packet, adapters, and analysis before inference. The breadth supplement freezes seven configurations and all 1,260 cells before inference on that known public challenge. Paper construction does not rerun any completed singleton study.

A post-r6 audit found provider hidden-thinking blocks in 275 source envelopes; scoring had used only final response text and parsed decisions. The reviewer derivative excludes those envelopes through a text-only allowlist and verifies source-to-derived hashes and scientific-payload equivalence. The anonymous bundle contains replayable item records, policy and prompt bindings, executable checks, formal checks, and privacy receipts. Its manifest derives the current file count from the generated tree. Fail-closed scans reject author identity, repository or workflow identifiers, provider response identifiers, hidden reasoning, and unknown record fields.

\section{AI Assistance Disclosure}

Generative AI tools assisted literature search, planning, code and benchmark drafting, debugging, analysis checks, and manuscript editing. The author chose the research question, defined the policy and protocol, made the study and analysis decisions, verified the benchmark, results, and references, and is responsible for the manuscript. Quantitative results come from validated study artifacts. AI systems are neither authors nor human annotators.

\section{Conclusion}

\cpb provides a replayable evaluation protocol for policy-constrained evidence-state transitions. The hidden r6 Qwen configuration reaches .8600/.7533 end-to-end item/pair accuracy at .9067/.8200 coverage, versus .9485/.9187 conditionally. On the action-balanced public challenge, Qwen and Flash reach .7333/.5333 and .5000/.1444 end to end, with conditional values of .8919/.7500 and .6164/.2131. Their full matrices show \demote{} recall of 4/30 and 0/30. The artifact demonstrates a failure-aware evaluation protocol; it does not establish broad benchmark validity or general research ability.

The seven-configuration breadth supplement retains all 1,260 terminal cells but obtains only 78 valid decisions. End-to-end item/pair accuracy ranges from .0000--.2167/.0000--.0444; four routes have undefined conditional metrics, and the others use no more than 65 valid items. Complete matrices and action recalls show zero \demote{} recall for every route. These are service-configuration outcomes on a known public challenge, not a causal ablation or model-family ranking. The old Flash 401 lane remains a failure of its old frozen service interface; a later executable route cannot recover its efficacy estimate. The deterministic executor is a definitional policy ceiling, while Naive Bayes remains a weak baseline. Naturalistic cases, stronger learned baselines, repeated configurations, and independent policy assessment are the next discriminating tests.

\bibliographystyle{plain}
\bibliography{references}

\end{document}
